\documentclass[../../main.tex]{subfiles}% 注意这里的文件路径不能用 ./main.tex ,否则用latexmk编译子文件会报错
\graphicspath{{\subfix{./image/}}} % 指定图片目录,后续可以直接使用图片文件名
% 注意这里的文件路径不能用 ../../image/ ,否则用latexmk编译子文件会报错

% 例如:
% \begin{figure}[H]
% \centering
% \includegraphics[scale=0.3]{图.png}
% \caption{}
% \label{figure:图}
% \end{figure}
% 注意:上述\label{}一定要放在\caption{}之后,否则引用图片序号会只会显示??.

\begin{document}

\section{内积空间}

\subsection{定义与基本性质}

\begin{definition}
线性空间 $\mathscr{X}$ 上的一个二元函数 $a(\cdot,\cdot):\mathscr{X}\times\mathscr{X}\to\mathbb{K}$，称为是\textbf{共轭双线性函数}，如果
\begin{enumerate}[(1)]
\item $a(x,\alpha_1y_1+\alpha_2y_2)=\overline{\alpha_1}a(x,y_1)+\overline{\alpha_2}a(x,y_2)$,
\item $a(\alpha_1x_1+\alpha_2x_2,y)=\alpha_1a(x_1,y)+\alpha_2a(x_2,y)$,
\end{enumerate}
其中 $\forall x,y,x_1,x_2,y_1,y_2\in\mathscr{X},\forall\alpha_1,\alpha_2\in\mathbb{K}$. 称由
\begin{align*}
q(x)\triangleq a(x,x)\quad(\forall x\in\mathscr{X})
\end{align*}
定义的函数为 $\mathscr{X}$ 上由 $a$ 诱导的\textbf{二次型}.
\end{definition}

\begin{proposition}
设 $a$ 是 $\mathscr{X}$ 上的共轭双线性函数，$q$ 是由 $a$ 诱导的二次型，那么
\begin{align*}
q(x)\in\mathbb{R}\ (\forall x\in\mathscr{X})\iff a(x,y)=\overline{a(y,x)}\quad(\forall x,y\in\mathscr{X}).
\end{align*}
\end{proposition}
\begin{proof}
充分性是显然的,下证必要性. 从
\begin{align*}
q(x+y)=\overline{q(x+y)}\quad(\forall x,y\in\mathscr{X}),
\end{align*}
容易推出
\begin{align}
a(x,y)+a(y,x)=\overline{a(x,y)}+\overline{a(y,x)}.\label{eq:::--fiojoi3048HUAAIJOW1}
\end{align}
在 \eqref{eq:::--fiojoi3048HUAAIJOW1} 式中换 $y$ 为 $\mathrm{i}y$，即得
\begin{align}
-a(x,y)+a(y,x)=\overline{a(x,y)}-\overline{a(y,x)}.\label{eq:::--fiojoi3048HUAAIJOW2}
\end{align}
\eqref{eq:::--fiojoi3048HUAAIJOW1} 式与 \eqref{eq:::--fiojoi3048HUAAIJOW2} 式相减即得 $a(x,y)=\overline{a(y,x)}$.

\end{proof}

\begin{definition}
\begin{enumerate}
\item 线性空间 $\mathscr{X}$ 上的一个共轭双线性函数 $(\cdot,\cdot):\mathscr{X}\times\mathscr{X}\to\mathbb{K}$ 称为是一个\textbf{半内积}，如果它满足:
\begin{enumerate}[(1)]
\item 共轭对称性:$(x,y)=\overline{(y,x)}\quad(\forall x,y\in\mathscr{X})$;
\item 非负定性:$(x,x)\geqslant 0\ (\forall x\in\mathscr{X})$ .
\end{enumerate}
具有内积的线性空间称为\textbf{半内积空间}.

\item 线性空间 $\mathscr{X}$ 上的一个共轭双线性函数 $(\cdot,\cdot):\mathscr{X}\times\mathscr{X}\to\mathbb{K}$ 称为是一个\textbf{内积}，如果它满足:
\begin{enumerate}[(1)]
\item 共轭对称性:$(x,y)=\overline{(y,x)}\quad(\forall x,y\in\mathscr{X})$;
\item 正定性:$(x,x)\geqslant 0\quad(\forall x\in\mathscr{X})$, $(x,x)=0\iff x=\theta$ .
\end{enumerate}
具有内积的线性空间称为\textbf{内积空间}，记作 $(\mathscr{X},(\cdot,\cdot))$.
\end{enumerate}
\end{definition}

\begin{proposition}
\begin{enumerate}[(1)]
\item $\mathbb{R}^n,\mathbb{C}^n$ 都是内积空间，它们的内积分别定义为
\begin{align*}
(x,y)=\sum\limits_{i=1}^n x_iy_i\quad(\forall x,y\in\mathbb{R}^n),
\end{align*}
\begin{align*}
(x,y)=\sum\limits_{i=1}^n x_i\overline{y_i}\quad(\forall x,y\in\mathbb{C}^n),
\end{align*}
其中 $x=(x_1,x_2,\cdots,x_n),y=(y_1,y_2,\cdots,y_n)$.

\item $l^2$ 空间是内积空间(定义见\refexas{examples:常见$B$空间11111}注(2))，规定内积
\begin{align*}
(x,y)=\sum\limits_{i=1}^\infty x_i\overline{y_i},
\end{align*}
其中 $x=(x_1,x_2,\cdots,x_i,\cdots),y=(y_1,y_2,\cdots,y_i,\cdots)\in l^2$.

\item $L^2(\Omega,\mu)$ 是内积空间(定义见\refexas{examples:常见$B$空间11111})，规定内积
\begin{align*}
(u,v)=\int_\Omega u(x)\cdot\overline{v(x)}\mathrm{d}\mu,
\end{align*}
其中 $u,v\in L^2(\Omega,\mu)$.

\item 在空间 $C^k(\overline{\Omega})$ 中(定义见\refexas{examples:常见的$B$空间222222})，规定内积
\begin{align*}
(u,v)=\sum\limits_{|\alpha|\leqslant k}\int_\Omega \partial^\alpha u(x)\cdot\overline{\partial^\alpha v(x)}\mathrm{d}x\quad(\forall u,v\in C^k(\overline{\Omega})).
\end{align*}
那么 $(C^k(\overline{\Omega}),(\cdot,\cdot))$ 是一个内积空间.
\end{enumerate}
\end{proposition}

\begin{proposition}[Cauchy-Schwarz不等式]\label{proposition:泛函分析--Cauchy-Schwarz不等式}
设 $(\mathscr{X},(\cdot,\cdot))$ 是内积空间. 若令
\begin{align*}
\|x\|=(x,x)^\frac{1}{2}\quad(\forall x\in\mathscr{X}),
\end{align*}
则有
\begin{align*}
|(x,y)|\leqslant\|x\|\cdot\|y\|\quad(\forall x,y\in\mathscr{X}),
\end{align*}
而且上式中的等号成立当且仅当 $x$ 与 $y$ 线性相关时成立.
\end{proposition}

\begin{proposition}
设 $a$ 是线性空间 $\mathscr{X}$ 上的共轭双线性函数，$q(x)$ 是由 $a$ 诱导的二次型. 如果
\begin{align*}
q(x)\geqslant 0\quad(\forall x\in\mathscr{X})\quad\text{且}\quad q(x)=0\iff x=\theta,
\end{align*}
那么
\begin{align}
|a(x,y)|\leqslant [q(x)q(y)]^\frac{1}{2}\quad(\forall x,y\in\mathscr{X}),\label{eq:::--fiojoi3048HUAAIJOW5}
\end{align}
而且上式中的等号成立当且仅当 $x$ 与 $y$ 线性相关时成立.
\end{proposition}
\begin{proof}
不妨设 $y\neq\theta$, 对 $\forall\lambda\in\mathbb{K}$ 考察
\begin{align}
q(x+\lambda y)=q(x)+\overline{\lambda}a(x,y)+\lambda a(y,x)+|\lambda|^2q(y)\geqslant 0.\label{eq:::--fiojoi3048HUAAIJOW6}
\end{align}
取 $\lambda=-a(x,y)/q(y)$, 因为 $a(x,y)=\overline{a(y,x)}$ (这是由假设 $q(x)\geqslant 0$, 根据前面的命题推出的), 所以
\begin{align*}
q(x)-\frac{2|a(x,y)|^2}{q(y)}+\frac{|a(x,y)|^2}{q(y)}\geqslant 0,
\end{align*}
由此立得 \eqref{eq:::--fiojoi3048HUAAIJOW5} 式. 又当 $x=-\lambda y\ (\lambda\in\mathbb{K})$ 时, \eqref{eq:::--fiojoi3048HUAAIJOW5} 式中的等号成立. 反之, 若 \eqref{eq:::--fiojoi3048HUAAIJOW5} 式中的等号成立, 则 \eqref{eq:::--fiojoi3048HUAAIJOW6} 式中的等号成立, 从而 $x=-\lambda y\ (\lambda\in\mathbb{K})$.

\end{proof}

\begin{proposition}
内积空间 $(\mathscr{X},(\cdot,\cdot))$ 按
\begin{align}\label{eq:::--fiojoi3048HUAAIJOW3}
\|x\|=(x,x)^\frac{1}{2}\quad(\forall x\in\mathscr{X})
\end{align}
定义范数$\|\cdot\|$,是严格凸的 $B^*$ 空间.并且内积 $(x,y)$ 是 $\mathscr{X}\times\mathscr{X}$ 上关于范数 $\|\cdot\|$ 的连续函数.
\end{proposition}
\begin{proof}
只要证明按 \eqref{eq:::--fiojoi3048HUAAIJOW3} 式定义的 $\|\cdot\|$ 是范数. 事实上, 正定性和齐次性都是显然成立的, 下面来验证三角不等式.
\begin{align*}
\|x+y\|^2&=(x+y,x+y)
=(x,x)+(x,y)+(y,x)+(y,y)\\
&\leqslant\|x\|^2+2\|x\|\|y\|+\|y\|^2
=(\|x\|+\|y\|)^2\quad(\forall x,y\in\mathscr{X}).
\end{align*}

对$\forall 0<\lambda<1$, 根据\hyperref[proposition:泛函分析--Cauchy-Schwarz不等式]{Cauchy-Schwarz不等式}我们有
\begin{align*}
\|\lambda x+(1-\lambda)y\|^2
&=\lambda^2\|x\|^2+2\lambda(1-\lambda)\text{Re}(x,y)+(1-\lambda)^2\|y\|^2\\
&<[\lambda+(1-\lambda)]^2=1\quad(\text{当}\ \|x\|=\|y\|=1,x\neq y).
\end{align*}

再证明内积的连续性.设 $x_n\to x,y_n\to y$. 那么 $\|x_n\|$ 和 $\|y_n\|$ 有界, 用 $M$ 表示它们的一个上界, 便有
\begin{align*}
|(x_n,y_n)-(x,y)|
&\leqslant|(x_n,y_n)-(x,y_n)|+|(x,y_n)-(x,y)|\\
&\leqslant\|x_n-x\|\cdot\|y_n\|+\|x\|\cdot\|y_n-y\|\\
&\leqslant M\|x_n-x\|+\|x\|\|y_n-y\|\to 0\quad(n\to\infty).
\end{align*}

\end{proof}

\begin{proposition}\label{proposition:平行四边形等式--泛函分析}
在 $B^*$ 空间 $(\mathscr{X},\|\cdot\|)$ 中, 在 $\mathscr{X}$ 上可引入一个内积 $(\cdot,\cdot)$ 适合
\begin{align}
(x,x)^\frac{1}{2}=\|x\|\quad(\forall x\in\mathscr{X}),\label{eq:::--fiojoi3048HUAAIJOW7}
\end{align}
当且仅当范数 $\|\cdot\|$ 满足如下\textbf{平行四边形等式}:
\begin{align*}
\|x+y\|^2+\|x-y\|^2=2(\|x\|^2+\|y\|^2)\quad(\forall x,y\in\mathscr{X}).
\end{align*}
\end{proposition}

\begin{proof}
必要性可通过直接计算得到. 为了证充分性, 令
\begin{align*}
(x,y)=
\begin{cases}
\frac{1}{4}\left(\|x+y\|^2-\|x-y\|^2\right), & \text{当}\ \mathbb{K}=\mathbb{R}, \\
\frac{1}{4}\left(\|x+y\|^2-\|x-y\|^2+i\|x+iy\|^2-i\|x-iy\|^2\right), & \text{当}\ \mathbb{K}=\mathbb{C}.
\end{cases}
\end{align*}
容易验证它是一个满足 \eqref{eq:::--fiojoi3048HUAAIJOW7} 式的内积.

\end{proof}

\begin{definition}
完备的内积空间称为 $\mathbf{Hilbert}$\textbf{空间}.
\end{definition}

\begin{examples}
$\mathbb{R}^n,\mathbb{C}^n$, $l^2$ 空间, $L^2(\Omega,\mu)$ 都是 Hilbert 空间.
\end{examples}

\begin{lemma}[Poincaré不等式]\label{lemma:(庞加莱)Poincaré不等式}
设 $C_0^m(\Omega)$ 表示有界开区域 $\Omega\subset\mathbb{R}^n$ 上一切 $m$ 次连续可微, 并在边界 $\partial\Omega$ 的某邻域内为 $0$ 的函数集合, 即
\begin{align*}
C_0^m(\Omega)=\{u\in C^m(\overline{\Omega})\mid u(x)=0,\text{当}\ x\in\partial\Omega\ \text{的某邻域}\}.
\end{align*}
那么 $\forall u\in C_0^m(\Omega)$ 有
\begin{align}
\sum\limits_{|\alpha|<m}\int_\Omega |\partial^\alpha u(x)|^2\mathrm{d}x\leqslant C\sum\limits_{|\alpha|=m}\int_\Omega |\partial^\alpha u(x)|^2\mathrm{d}x,\label{eq:::--fiojoi3048HUAAIJOW9}
\end{align}
其中 $C$ 是仅依赖于区域 $\Omega$ 及 $m$ 的常数.
\end{lemma}
\begin{remark}
这个引理表明在 $C_0^m(\Omega)$ 上,
\begin{align}
\|u\|_m\triangleq\left(\sum\limits_{|\alpha|=m}\int_\Omega |\partial^\alpha u(x)|^2\mathrm{d}x\right)^\frac{1}{2}\label{eq:::--fiojoi3048HUAAIJOW12}
\end{align}
和
\begin{align*}
\|u\|\triangleq\left(\sum\limits_{|\alpha|\leqslant m}\int_\Omega |\partial^\alpha u(x)|^2\mathrm{d}x\right)^\frac{1}{2}
\end{align*}
是一对等价范数. 记 $C_0^m(\Omega)$ 按 \eqref{eq:::--fiojoi3048HUAAIJOW12} 式完备化后的空间为 $H_0^m(\Omega)$. 它是 $H^m(\Omega)$ 的一个闭子空间.
\end{remark}
\begin{proof}
因为 $\Omega$ 是有界的, 我们可以把 $\Omega$ 放在某个边长为 $a$ 的立方体 $\Omega_1$ 内, 适当选择坐标系, 使得
\begin{align*}
\Omega_1=\{(x_1,x_2,\cdots,x_n)\in\mathbb{R}^n\mid 0\leqslant x_i\leqslant a\ (i=1,2,\cdots,n)\}.
\end{align*}
在 $\Omega_1\setminus\Omega$ 上补充定义 $u=0$, 经补充定义后, $u(x)$ 在 $\Omega_1$ 上 $m$ 次连续可微, 而且在边界上等于 $0$. $\forall x\in\Omega_1$,
\begin{align*}
u(x)=\int_0^{x_1}\frac{\partial u}{\partial t}(t,x_2,\cdots,x_n)\mathrm{d}t.
\end{align*}
再利用\hyperref[proposition:泛函分析--Cauchy-Schwarz不等式]{Cauchy-Schwarz不等式}, 我们有
\begin{align}
|u(x)|^2\leqslant a\int_0^a\left|\frac{\partial u}{\partial x_1}\right|^2\mathrm{d}x_1.\label{eq:::--fiojoi3048HUAAIJOW10}
\end{align}
在 $\Omega_1$ 上对不等式 \eqref{eq:::--fiojoi3048HUAAIJOW10}积分得
\begin{align}\label{eq:::--fiojoi3048HUAAIJOW11}
\int_\Omega |u(x)|^2\mathrm{d}x
\leqslant a^2\int_\Omega \left|\frac{\partial u}{\partial x_1}\right|^2\mathrm{d}x
\leqslant a^2\int_\Omega |\text{grad}\ u(x)|^2\mathrm{d}x.
\end{align}
然后逐次应用不等式 \eqref{eq:::--fiojoi3048HUAAIJOW11} 于 $\partial^\alpha u(x)\ (|\alpha|<m)$, 即得不等式 \eqref{eq:::--fiojoi3048HUAAIJOW9}.

\end{proof}

\begin{example}
$H_0^m(\Omega)$ 是一个 Hilbert 空间, 其内积定义为
\begin{align*}
(u,v)_m=\sum\limits_{|\alpha|=m}\int_\Omega \partial^\alpha u(x)\cdot\overline{\partial^\alpha v(x)}\mathrm{d}x\quad(\forall u,v\in C_0^m(\Omega)).
\end{align*}
\end{example}

\begin{remark}
当 $\Omega$ 的边界 $\partial\Omega$ 具有光滑的法向导数时, $H_0^m(\Omega)$ 中的元素实际上是满足边界条件:
\begin{align*}
u|_{\partial\Omega}=\frac{\partial u}{\partial n}\bigg|_{\partial\Omega}=\cdots=\left(\frac{\partial}{\partial n}\right)^{m-1}u\bigg|_{\partial\Omega}=0
\end{align*}
的 $C^m(\Omega)$ 函数 $u$ 的一种推广, 其中 $\frac{\partial}{\partial n}$ 是 $\partial\Omega$ 上的法向导数.
\end{remark}


\subsection{正交与正交基}

\begin{definition}
内积空间 $\mathscr{X}$ 上的两个元素 $x$ 与 $y$ 称为是\textbf{正交的}, 是指
\begin{align*}
(x,y)=0,
\end{align*}
记作 $x\perp y$. 又设 $M$ 是 $\mathscr{X}$ 的一个非空子集, $x\in\mathscr{X}$. 若对 $\forall y\in M$ 都有 $x\perp y$, 则称 $x$ 与 $M$ \textbf{正交}, 记作 $x\perp M$. 此外我们还称集合
\begin{align*}
\{x\in\mathscr{X}\mid x\perp M\}
\end{align*}
为 $M$ 的\textbf{正交补}, 记作 $M^\perp$.
\end{definition}

\begin{proposition}
设 $\mathscr{X}$ 是内积空间, $M$ 是 $\mathscr{X}$ 的一个非空子集.
\begin{enumerate}[(1)]
\item 若 $x\perp y_i\ (i=1,2)$, 则 $x\perp\lambda_1y_1+\lambda_2y_2\ (\forall\lambda_1,\lambda_2\in\mathbb{K})$.
\item 若 $x=y+z$, 且 $y\perp z$, 则 $\|x\|^2=\|y\|^2+\|z\|^2$.
\item 若 $x\perp y_n\ (n\in\mathbb{N})$, 且 $y_n\to y$, 则 $x\perp y$.
\item 若 $x\perp M$, 则 $x\perp\text{span}M$.
\item $M^\perp$ 是 $\mathscr{X}$ 的一个闭线性子空间.
\end{enumerate}
\end{proposition}

\begin{definition}
设 $\mathscr{X}$ 是一个内积空间,$\theta$是零向量, 集合 $S=\{e_\alpha\mid\alpha\in A\}$ 是 $\mathscr{X}$ 的一个子集. 称 $S$ 为\textbf{正交集}, 是指
\begin{align*}
e_\alpha\perp e_\beta\quad(\text{当}\ \alpha\neq\beta,\forall\alpha,\beta\in A).
\end{align*}
如果还有 $\|e_\alpha\|=1\ (\forall\alpha\in A)$, 则称 $S$ 为\textbf{正交规范集}. 又如果在 $\mathscr{X}$ 中不存在非零元与 $S$ 正交, 即 $S^\perp=\{\theta\}$, 那么称 $S$ 为\textbf{完备的}.
\end{definition}

\begin{lemma}[Zorn引理]\label{lemma:泛函分析--Zorn引理}
设 $\mathscr{X}$ 是一个半序集. 如果它的每一个全序子集有一个上界, 那么 $\mathscr{X}$ 有一个极大元.
\end{lemma}

\begin{proposition}
非 $\{\theta\}$ 内积空间 $\mathscr{X}$ 中必存在完备正交集.
\end{proposition}
\begin{proof}
因为 $\mathscr{X}\neq\{\theta\}$, 所以 $\mathscr{X}$ 中的正交集依包含关系构成一个半序集类, 并且每个全序子集类有一个上界, 就是这些集之并集. 依\hyperref[lemma:泛函分析--Zorn引理]{Zorn 引理}, 这个半序集类有极大元. 我们来证明: 这个极大元 (记作 $S$) 就是完备正交集. 因若不然, 则必 $\exists x_0\in S^\perp,x_0\neq\theta$, 令 $S_1=\{x_0\}\cup S$, 得到 $S_1$ 还是正交集, 并且 $S\subsetneq S_1$, 这便与 $S$ 的极大性相矛盾!

\end{proof}

\begin{definition}
内积空间 $\mathscr{X}$ 中的正交规范集 $S=\{e_\alpha\mid\alpha\in A\}$, 称为一个\textbf{基} 或\textbf{封闭的} 是指 $\forall x\in\mathscr{X}$, 有下列表示:
\begin{align*}
x=\sum\limits_{\alpha\in A}(x,e_\alpha)e_\alpha,
\end{align*}
其中 $\{(x,e_\alpha)\mid\alpha\in A\}$ 称为 $x$ 关于基 $\{e_\alpha\mid\alpha\in A\}$ 的 $\mathbf{Fourier}$ \textbf{系数}.
\end{definition}

\begin{theorem}[Bessel不等式]\label{theorem:泛函分析---Bessel不等式}
设 $\mathscr{X}$ 是一个内积空间. 如果 $S=\{e_\alpha\mid\alpha\in A\}$ 是 $\mathscr{X}$ 中的正交规范集, 那么 $\forall x\in\mathscr{X}$, 有
\begin{align}
\sum\limits_{\alpha\in A}|(x,e_\alpha)|^2\leqslant\|x\|^2.\label{eq:::--fiojoi3048HUAAIJOW16}
\end{align}
\end{theorem}
\begin{proof}
首先对 $A$ 的任意有限子集, 不妨设它们是 $1,2,\cdots,n$, 证明
\begin{align}
\sum\limits_{i=1}^n|(x,e_i)|^2\leqslant\|x\|^2.\label{eq:::--fiojoi3048HUAAIJOW17}
\end{align}
因为
\begin{align*}
0\leqslant\left\|x-\sum\limits_{i=1}^n(x,e_i)e_i\right\|^2
=\left(x-\sum\limits_{i=1}^n(x,e_i)e_i,x-\sum\limits_{j=1}^n(x,e_j)e_j\right)=\|x\|^2-\sum\limits_{i=1}^n|(x,e_i)|^2,
\end{align*}
所以 \eqref{eq:::--fiojoi3048HUAAIJOW17} 式成立. 由此可见, 对 $\forall n\in\mathbb{N}$, 适合 $|(x,e_\alpha)|>1/n$ 的 $\alpha\in A$ 至多只有有限多个, 从而 $(x,e_\alpha)\neq 0$ 的 $\alpha\in A$ 至多有可数多个. 于是 \eqref{eq:::--fiojoi3048HUAAIJOW16} 式的左端实际上是至多可数项求和的级数. 再由 \eqref{eq:::--fiojoi3048HUAAIJOW17} 式我们有
\begin{align*}
\sum\limits_{\alpha\in A_f}|(x,e_\alpha)|^2\leqslant\|x\|^2,
\end{align*}
其中 $A_f$ 表示 $A$ 的任意有限子集. 由此立得 \eqref{eq:::--fiojoi3048HUAAIJOW16} 式.

\end{proof}

\begin{corollary}
假设 $\mathscr{X}$ 是 Hilbert 空间, 且 $\{e_\alpha\mid\alpha\in A\}$ 是 $\mathscr{X}$ 中的正交规范集. 那么对 $\forall x\in\mathscr{X}$, 有
\begin{align*}
\sum\limits_{\alpha\in A}(x,e_\alpha)e_\alpha\in\mathscr{X},
\end{align*}
且
\begin{align}
\left\|x-\sum\limits_{\alpha\in A}(x,e_\alpha)e_\alpha\right\|^2=\|x\|^2-\sum\limits_{\alpha\in A}|(x,e_\alpha)|^2.\label{eq:::--fiojoi3048HUAAIJOW18}
\end{align}
\end{corollary}

\begin{proof}
不妨设使得 $(x,e_\alpha)\neq 0$ 的可数多个 $\alpha\in A$ 是 $1,2,\cdots,n,\cdots$, 那么
\begin{align*}
\sum\limits_{\alpha\in A}(x,e_\alpha)e_\alpha=\sum\limits_{n=1}^\infty(x,e_n)e_n,
\end{align*}
并由\hyperref[theorem:泛函分析---Bessel不等式]{Bessel不等式}可知 $\sum\limits_{n=1}^\infty|(x,e_n)|^2$ 收敛, 因此有
\begin{align*}
\left\|\sum\limits_{n=m}^{m+p}(x,e_n)e_n\right\|^2=\sum\limits_{n=m}^{m+p}|(x,e_n)|^2\to 0\quad(m\to\infty,\forall p\in\mathbb{N}).
\end{align*}
于是, $\left\{x_m=\sum\limits_{n=1}^m(x,e_n)e_n\right\}$ 是基本列, 从而
\begin{align*}
\sum\limits_{\alpha\in A}(x,e_\alpha)e_\alpha=\sum\limits_{n=1}^\infty(x,e_n)e_n=\lim_{m\to\infty}x_m\in\mathscr{X}.
\end{align*}
又因为 $x-\sum\limits_{n=1}^\infty(x,e_n)e_n\perp\sum\limits_{n=1}^\infty(x,e_n)e_n$, 所以
\begin{align*}
\left\|x-\sum\limits_{n=1}^\infty(x,e_n)e_n\right\|^2=\|x\|^2-\sum\limits_{n=1}^\infty|(x,e_n)|^2.
\end{align*}
这就是 \eqref{eq:::--fiojoi3048HUAAIJOW18} 式.

\end{proof}

\begin{theorem}\label{theorem:泛函分析--Parseval恒等式}
设$\mathscr{X}$是一个Hilbert空间，若$S=\{e_\alpha \mid \alpha \in A\}$是$\mathscr{X}$中的正交规范集，则如下三条等价：
\begin{enumerate}[(1)]
\item $S$是封闭的；
\item $S$是完备的；
\item $\mathbf{Parseval}$\textbf{等式}
\begin{align}
\|x\|^2 = \sum\limits_{\alpha \in A} |(x, e_\alpha)|^2 \quad (\forall x \in \mathscr{X}). \label{eq:parseval}
\end{align}
成立。
\end{enumerate}
\end{theorem}
\begin{proof}
(1)$\implies$(2). 若$S$不完备，则$\exists x \in \mathscr{X} \setminus \{\theta\}$，使得
\begin{align*}
(x, e_\alpha) = 0 \quad (\forall \alpha \in A).
\end{align*}
但由封闭性有$x = \sum\limits_{\alpha \in A} (x, e_\alpha) e_\alpha = \theta$，矛盾。

(2)$\implies$(3). 若$\exists x \in \mathscr{X}$使Parseval等式\eqref{eq:parseval}不成立，则由
\begin{align*}
\left\|x - \sum\limits_{\alpha \in A} (x, e_\alpha) e_\alpha\right\|^2 = \|x\|^2 - \sum\limits_{\alpha \in A} |(x, e_\alpha)|^2 > 0.
\end{align*}
于是$y \triangleq x - \sum\limits_{\alpha \in A} (x, e_\alpha) e_\alpha \neq \theta$，但$y \in S^\perp$。这与$S$完备性矛盾。

(3)$\implies$(1). 联合Parseval等式\eqref{eq:parseval}与
\begin{align*}
\left\|x - \sum\limits_{\alpha \in A} (x, e_\alpha) e_\alpha\right\|^2 = \|x\|^2 - \sum\limits_{\alpha \in A} |(x, e_\alpha)|^2 = 0.
\end{align*}
因此有
\begin{align*}
x = \sum\limits_{\alpha \in A} (x, e_\alpha) e_\alpha.
\end{align*}

\end{proof}

\begin{example}
在$L^2[0, 2\pi]$上，
\begin{align*}
e_n(t) = \frac{1}{\sqrt{2\pi}} e^{\mathrm{i}nt} \quad (n=0, \pm1, \pm2, \cdots)
\end{align*}
是一组正交规范基。$\forall u \in L^2[0, 2\pi]$，对应的Fourier系数是
\begin{align*}
(u, e_n) = \frac{1}{\sqrt{2\pi}} \int_{0}^{2\pi} u(t) e^{-\mathrm{i}nt} \mathrm{d}t \quad (n=0, \pm1, \pm2, \cdots).
\end{align*}
\end{example}

\begin{example}
在$l^2$空间上，
\begin{align*}
e_n = (\underbrace{0, 0, \cdots, 0, 1}_{n}, 0, \cdots) \quad (n=1,2,3,\cdots)
\end{align*}
是一组正交规范基。
\end{example}

\begin{proposition}\label{proposition:H2D空间--泛函分析}
设$D$是$\mathbb{C}$中的单位开圆域，$H^2(D)$表示在$D$内满足
\begin{align*}
\iint_{D} |u(z)|^2 \mathrm{d}x\mathrm{d}y < \infty \quad (z=x+\mathrm{i}y)
\end{align*}
的解析函数全体组成的空间。规定内积为
\begin{align*}
(u, v) = \iint_{D} u(z) \overline{v(z)} \mathrm{d}x\mathrm{d}y.
\end{align*}
这时函数组
\begin{align*}
\varphi_n(z) = \sqrt{\frac{n}{\pi}} z^{n-1} \quad (n=1,2,3,\cdots)
\end{align*}
是一组正交规范基。设$u(z) = \sum\limits_{k=0}^{\infty} b_k z^k \in H^2(D)$，它对应的Fourier系数是
\begin{align*}
(u, \varphi_n) = \iint_{D} \sum\limits_{k=0}^{\infty} b_k z^k \sqrt{\frac{n}{\pi}} \overline{z}^{n-1} \mathrm{d}x\mathrm{d}y 
= \sum\limits_{k=0}^{\infty} b_k \sqrt{\frac{\pi}{k+1}} (\varphi_{k+1}, \varphi_n) = b_{n-1} \sqrt{\frac{\pi}{n}} \quad (n=1,2,3,\cdots).
\end{align*}
\end{proposition}


\subsection{正交化和Hilbert空间的同构}

\begin{theorem}[Gram-Schmidt 正交化]\label{theorem:Gram-Schmidt 正交化}
设 $\mathscr{X}$ 是内积空间，$\{x_n\}_{n=1}^{\infty} \subset \mathscr{X}$ 是线性无关的向量列。则存在正交规范列 $\{y_n\}_{n=1}^{\infty} \subset \mathscr{X}$，使得对任意 $n \in \mathbb{N}$ 有
\[
\operatorname{span}\{y_1,\cdots,y_n\} = \operatorname{span}\{x_1,\cdots,x_n\}.
\]
且满足
\[
\begin{cases}
y_1 = x_1, \quad e_1 = \dfrac{y_1}{\|y_1\|},\\[6pt]
y_n = x_n - \displaystyle\sum_{i=1}^{n-1} \langle x_n, e_i \rangle e_i, \quad e_n = \dfrac{y_n}{\|y_n\|}, \quad n \ge 2.
\end{cases}
\]
这里 $\langle\cdot,\cdot\rangle$ 与 $\|\cdot\|$ 分别为 $\mathscr{X}$ 上的内积和诱导范数。

若 $\{x_n\}$ 张成 $\mathscr{X}$ 的某个子空间 $M$，则 $\{y_n\}$ 构成 $M$ 的一组标准正交基。
\end{theorem}

\begin{definition}
设$(\mathscr{X}_1, (\cdot, \cdot)_1)$与$(\mathscr{X}_2, (\cdot, \cdot)_2)$是两个内积空间，如果存在$\mathscr{X}_1 \to \mathscr{X}_2$的一个线性同构$T$，满足
\begin{align*}
(Tx, Ty)_2 = (x, y)_1 \quad (\forall x, y \in \mathscr{X}_1),
\end{align*}
则称内积空间$\mathscr{X}_1$与$\mathscr{X}_2$是\textbf{同构的}。
\end{definition}

\begin{theorem}
Hilbert空间$\mathscr{X}$是可分的当且仅当它有至多可数的正交规范基$S$。又若$S$的元素个数$N < \infty$，则$\mathscr{X}$同构于$\mathbb{K}^N$；若$N = \infty$，则$\mathscr{X}$同构于$l^2$。
\end{theorem}
\begin{proof}
{\heiti 必要性:} 设$\{x_n\}_{n=1}^{\infty}$是$\mathscr{X}$中的可数稠密子集，那么其中必存在一个线性无关的子集$\{y_n\}_{n=1}^{N}$（$N < \infty$或$N = \infty$），使得
\begin{align*}
\overline{\operatorname{span}\{y_n\}_{n=1}^{N}} = \overline{\operatorname{span}\{x_n\}_{n=1}^{\infty}} = \mathscr{X}.
\end{align*}
再对$\{y_n\}_{n=1}^{N}$应用Gram-Schmidt过程，便构造出一个正交规范集$\{e_n\}_{n=1}^{N}$。又因为
\begin{align*}
\overline{\operatorname{span}\{e_n\}_{n=1}^{N}} = \overline{\operatorname{span}\{y_n\}_{n=1}^{N}} = \mathscr{X},
\end{align*}
所以$\{e_n\}_{n=1}^{N}$是正交规范基。

{\heiti 充分性} 设$\{e_n\}_{n=1}^{N}$（$N < \infty$或$N = \infty$）是$\mathscr{X}$的正交规范基，那么集合
\begin{align*}
\left\{x = \sum\limits_{n=1}^{N} a_n e_n \mid \operatorname{Re}a_n \text{ 与 } \operatorname{Im}a_n \text{ 皆为有理数}\right\}
\end{align*}
是$\mathscr{X}$中的稠密子集。从而$\mathscr{X}$是可分的。

对于正交规范基$\{e_n\}_{n=1}^{N}$（$N < \infty$或$N = \infty$），做对应
\begin{align*}
T: x \mapsto \{(x, e_n)\}_{n=1}^{N} \quad (\forall x \in \mathscr{X}).
\end{align*}
根据\hyperref[theorem:泛函分析--Parseval恒等式]{Parseval等式}我们有
\begin{align*}
\|x\|^2 = \sum\limits_{n=1}^{N} |(x, e_n)|^2 \quad (\forall x \in \mathscr{X}).
\end{align*}
由此可见，对应$T$是$\mathscr{X} \to \mathbb{K}^N$（当$N < \infty$）或$\mathscr{X} \to l^2$（当$N = \infty$）的一对一在上线性同构。此外，
\begin{align*}
(x, y) = \left( \sum\limits_{i=1}^{N} (x, e_i) e_i, \sum\limits_{j=1}^{N} (y, e_j) e_j \right) 
= \sum\limits_{i=1}^{N} (x, e_i) \overline{(y, e_i)} \quad (\forall x, y \in \mathscr{X}).
\end{align*}
因此$T$还是保持内积的。于是当$N < \infty$时，$\mathscr{X}$同构于$\mathbb{K}^N$；而当$N = \infty$时，$\mathscr{X}$同构于$l^2$。

\end{proof}


\subsection{再论最佳逼近问题}

\begin{theorem}\label{theorem:泛函分析---定理1.1611}
如果$C$是Hilbert空间$\mathscr{X}$中的一个闭凸子集，那么在$C$上存在唯一元素$x_0$取到最小范数。
\end{theorem}
\begin{proof}
存在性. 若零向量$\theta \in C$，则$x_0 = \theta$；若零向量$\theta \notin C$，则
\begin{align*}
d \triangleq \inf\limits_{z \in C} \|z\| > 0.
\end{align*}
由下确界定义，$\forall n \in \mathbb{N}$，$\exists x_n \in C$，使得
\begin{align}
d \leqslant \|x_n\| \leqslant d + \frac{1}{n} \quad (n=1,2,\cdots). \label{eq:min_norm_seq}
\end{align}
如果$\{x_n\}$有极限$x_0$，那么由$C$的闭性，$x_0 \in C$，并由\eqref{eq:min_norm_seq}式，$\|x_0\| = d$，即$x_0$取到了$C$上元素的最小范数。为了$\{x_n\}$有极限，只需验证$\{x_n\}$是一个基本列。这要用到\hyperref[proposition:平行四边形等式--泛函分析]{平行四边形等式}:
\begin{align*}
\|x_m - x_n\|^2 = 2(\|x_m\|^2 + \|x_n\|^2) - 4\left\|\frac{x_m + x_n}{2}\right\|^2 
\leqslant 2\left[ \left(d + \frac{1}{n}\right)^2 + \left(d + \frac{1}{m}\right)^2 \right] - 4d^2 \to 0 \quad (\text{当 } n, m \to \infty).
\end{align*}

唯一性. 如果有$x_0, \hat{x}_0 \in C$使得$\|x_0\| = \|\hat{x}_0\| = d$，那么
\begin{align*}
\|x_0 - \hat{x}_0\|^2 = 2(\|x_0\|^2 + \|\hat{x}_0\|^2) - 4\left\|\frac{x_0 + \hat{x}_0}{2}\right\|^2 
\leqslant 4d^2 - 4d^2 = 0,
\end{align*}
即得$x_0 = \hat{x}_0$。

\begin{figure}[H]
\centering
\includegraphics[scale=0.3]{泛函分析图1.6.1.png}
\caption{}
\label{figure:min_norm}
\end{figure}

\end{proof}

\begin{corollary}
若$C$是Hilbert空间$\mathscr{X}$中的一个闭凸子集，则对$\forall y \in \mathscr{X}$，$\exists x_0 \in C$，使得
\begin{align}
\|y - x_0\| = \inf\limits_{x \in C} \|x - y\|. \label{eq:best_approx}
\end{align}
\end{corollary}

\begin{proof}
只需考察集合$C - \{y\} \triangleq \{x - y \mid x \in C\}$，它显然还是$\mathscr{X}$上的一个闭凸子集。根据\refthe{theorem:泛函分析---定理1.1611}，$\exists z_0 \in C - \{y\}$，使得
\begin{align*}
\|z_0\| = \inf\limits_{z \in C - \{y\}} \|z\|.
\end{align*}
令$x_0 \triangleq z_0 + y$，便得到\eqref{eq:best_approx}式。

\end{proof}

\begin{corollary}
若$M$是Hilbert空间$\mathscr{X}$上的一个闭线性子空间，则对$\forall y \in \mathscr{X}$，$\exists x_0 \in M$，使得
\begin{align*}
\|y - x_0\| = \inf\limits_{x \in M} \|x - y\|.
\end{align*}
\end{corollary}

\begin{theorem}[最佳逼近元]\label{theorem:泛函分析--最佳逼近元的刻画}
设$C$是内积空间$\mathscr{X}$中的一个闭凸子集，$\forall y \in \mathscr{X}$，则$x_0$是$y$在$C$上的最佳逼近元当且仅当它适合
\begin{align}
\operatorname{Re}(y - x_0, x_0 - x) \geqslant 0 \quad (\forall x \in C). \label{eq:best_approx_char}
\end{align}
\end{theorem}
\begin{proof}
对$\forall x \in C$，考察函数
\begin{align*}
\varphi_x(t) = \|y - tx - (1-t)x_0\|^2 \quad (t \in [0,1]).
\end{align*}
显然，为了$x_0$是$y$在$C$上的最佳逼近元，当且仅当
\begin{align}
\varphi_x(t) \geqslant \varphi_x(0) \quad (\forall x \in C, \forall t \in [0,1]). \label{eq:phi_min}
\end{align}

下证\eqref{eq:best_approx_char}式成立$\iff$ \eqref{eq:phi_min}式成立。
因为
\begin{align*}
\varphi_x(t) = \|(y - x_0) + t(x_0 - x)\|^2 
= \|y - x_0\|^2 + 2t\operatorname{Re}(y - x_0, x_0 - x) + t^2\|x_0 - x\|^2,
\end{align*}
所以
\begin{align}
\varphi_x'(0) = 2\operatorname{Re}(y - x_0, x_0 - x). \label{eq:phi_deriv}
\end{align}
因此
\begin{align}
\eqref{eq:best_approx_char}式成立 \iff \varphi_x'(0) \geqslant 0. \label{eq:char_iff_deriv}
\end{align}
又因为
\begin{align*}
\varphi_x(t) - \varphi_x(0) = \varphi_x'(0)t + \|x_0 - x\|^2 t^2,
\end{align*}
所以
\begin{align}
\varphi_x'(0) \geqslant 0 \iff \eqref{eq:phi_min}式成立. \label{eq:deriv_iff_min}
\end{align}
联合\eqref{eq:char_iff_deriv}式与\eqref{eq:deriv_iff_min}式即得结论。

\end{proof}

\begin{corollary}
设$M$是Hilbert空间$\mathscr{X}$的一个闭线性子流形。对$\forall x \in \mathscr{X}$，$y$是$x$在$M$上的最佳逼近元当且仅当它适合
\begin{align}
x - y \perp M - \{y\} \triangleq \{w = z - y \mid \forall z \in M\}. \label{eq:best_approx_linear}
\end{align}
\end{corollary}
\begin{proof}
由\refthe{theorem:泛函分析--最佳逼近元的刻画}，$y$是$x$在$M$上的最佳逼近元当且仅当
\begin{align}
\operatorname{Re}(x - y, y - z) \geqslant 0 \quad (\forall z \in M). \label{eq:linear_char}
\end{align}
因为$M$是线性流形，所以$\forall z \in M$可表示为
\begin{align}
z = y + w \quad (w \in M - \{y\}). \label{eq:linear_rep}
\end{align}
注意到$M - \{y\}$是线性子空间，且当$z$跑遍$M$时，$w$跑遍$M - \{y\}$。将\eqref{eq:linear_rep}式代入\eqref{eq:linear_char}式得
\begin{align}
\operatorname{Re}(x - y, w) \leqslant 0 \quad (\forall w \in M - \{y\}). \label{eq:re_ineq}
\end{align}
在\eqref{eq:re_ineq}式中，用$-w$代替$w$，便推出
\begin{align}
\operatorname{Re}(x - y, w) = 0 \quad (\forall w \in M - \{y\}). \label{eq:re_eq}
\end{align}
进一步在\eqref{eq:re_eq}式中用$\mathrm{i}w$代替$w$，便推出
\begin{align*}
(x - y, w) = 0 \quad (\forall w \in M - \{y\}).
\end{align*}
这就是\eqref{eq:best_approx_linear}式。

\end{proof}

\begin{corollary}\label{corollary:泛函分析--推论289293}
设$M$是Hilbert空间$\mathscr{X}$的一个闭线性子空间，$M - \{y\} = M$,则$y$是$x$在$M$上的最佳逼近元当且仅当$x - y \perp M$。
\end{corollary}

\begin{corollary}[正交分解]\label{corollary:泛函分析--正交分解定理}
设$M$是Hilbert空间$\mathscr{X}$上的一个闭线性子空间。则对$\forall x \in \mathscr{X}$，存在下列唯一的\textbf{正交分解}：
\begin{align}
x = y + z \quad (y \in M, z \in M^\perp). \label{eq:orthogonal_decomp}
\end{align}
\end{corollary}
\begin{remark}
由$x$的正交分解产生的$y$称为$x$在$M$上的正交投影(见\reffig{figure:orthogonal_decomp})。
\begin{figure}[H]
\centering
\includegraphics[scale=0.3]{泛函分析--图1.6.5.png}
\caption{}
\label{figure:orthogonal_decomp}
\end{figure}
\end{remark}
\begin{proof}
取$y$为$x$在$M$上的最佳逼近元，而$z = x - y$，即为满足\eqref{eq:orthogonal_decomp}式的分解。又若还存在另外一种分解：
\begin{align*}
x = y' + z' \quad (y' \in M, z' \in M^\perp),
\end{align*}
则
\begin{align*}
y - y' = z - z' \in (M \cap M^\perp) = \{\theta\},
\end{align*}
即得分解的唯一性。

\end{proof}


\subsection{应用:最小二乘法}

\begin{example}\label{example:泛函分析--最小二乘法}
\begin{enumerate}[(1)]
\item \textbf{实际观测问题}. 已知量 $y$ 与量 $x_1, x_2, \cdots, x_n$ 之间呈线性关系:
\begin{align*}
y = \lambda_1 x_1 + \lambda_2 x_2 + \cdots + \lambda_n x_n.
\end{align*}
为确定未知系数 $\lambda_1, \lambda_2, \cdots, \lambda_n$，观测 $m$ 组数据（$m > n$），即
\begin{table}[H]
\centering
\caption{}
\begin{tabular}{ccccc}
$y^{(1)},$ & $x_1^{(1)},$ & $x_2^{(1)},$ & $\cdots,$ & $x_n^{(1)}$ \\
$y^{(2)},$ & $x_1^{(2)},$ & $x_2^{(2)},$ & $\cdots,$ & $x_n^{(2)}$ \\
$y^{(3)},$ & $x_1^{(3)},$ & $x_2^{(3)},$ & $\cdots,$ & $x_n^{(3)}$ \\
$\vdots$   & $\vdots$     & $\vdots$     & $\ddots$  & $\vdots$   \\
$y^{(m)},$ & $x_1^{(m)},$ & $x_2^{(m)},$ & $\cdots,$ & $x_n^{(m)}$ \\
\end{tabular}
\end{table}
求系数 $\lambda_1, \lambda_2, \cdots, \lambda_n$，使得
\begin{align*}
\min_{\alpha_1, \alpha_2, \cdots, \alpha_n} \sum\limits_{j=1}^{m} \left| y^{(j)} - \sum\limits_{i=1}^{n} \alpha_i x_i^{(j)} \right|^2 = \sum\limits_{j=1}^{m} \left| y^{(j)} - \sum\limits_{i=1}^{n} \lambda_i x_i^{(j)} \right|^2.
\end{align*}
该问题可视为在空间 $\mathbb{R}^m$ 中，求由数据表中后 $n$ 列向量张成的子空间上的元素，对第一列向量的最佳逼近.

\item \textbf{平方平均逼近}. 在函数逼近论中，对 $f \in L^2[a,b]$，用给定的 $n$ 个 $L^2[a,b]$ 函数 $\varphi_1, \varphi_2, \cdots, \varphi_n$ 的线性组合按平方平均意义求最佳逼近，即求系数 $\lambda_1, \lambda_2, \cdots, \lambda_n$，使得
\begin{align*}
\min_{\alpha_1, \alpha_2, \cdots, \alpha_n} \int_{a}^{b} \left| f(x) - \sum\limits_{i=1}^{n} \alpha_i \varphi_i(x) \right|^2 \mathrm{d}x = \int_{a}^{b} \left| f(x) - \sum\limits_{i=1}^{n} \lambda_i \varphi_i(x) \right|^2 \mathrm{d}x.
\end{align*}

\item \textbf{最佳估计问题}. 设 $(\Omega, \mathscr{B}, P)$ 是概率空间，$X, X_1, X_2, \cdots, X_n \in L^2(\Omega, \mathscr{B}, P)$，用 $X_1, X_2, \cdots, X_n$ 的线性组合估计 $X$，求系数 $\lambda_1, \lambda_2, \cdots, \lambda_n$，使得
\begin{align*}
\min_{(\alpha_1, \alpha_2, \cdots, \alpha_n) \in \mathbb{R}^n} \int_{\Omega} \left| X(\omega) - \sum\limits_{i=1}^{n} \alpha_i X_i(\omega) \right|^2 \mathrm{d}P(\omega) = \int_{\Omega} \left| X(\omega) - \sum\limits_{i=1}^{n} \lambda_i X_i(\omega) \right|^2 \mathrm{d}P(\omega).
\end{align*}
\end{enumerate}
\end{example}

\begin{theorem}
设 $\mathscr{X}$ 为 Hilbert 空间，$x_1, x_2, \cdots, x_n \in \mathscr{X}$ 线性无关，记 $\mathscr{M} = \operatorname{span}\{x_1, x_2, \cdots, x_n\}$，则存在唯一的 $x_0 = \sum\limits_{i=1}^{n} \lambda_i x_i \in \mathscr{M}$，使得
\begin{align*}
\left\| x - x_0 \right\| = \inf_{y \in \mathscr{M}} \left\| x - y \right\|,
\end{align*}
且
\begin{align*}
\lambda_i = \frac{
\begin{vmatrix}
(x_1, x_1) & \cdots & (x, x_1) & \cdots & (x_n, x_1) \\
(x_1, x_2) & \cdots & (x, x_2) & \cdots & (x_n, x_2) \\
\vdots & & \vdots & & \vdots \\
(x_1, x_n) & \cdots & (x, x_n) & \cdots & (x_n, x_n)
\end{vmatrix}
}{
\begin{vmatrix}
(x_1, x_1) & \cdots & (x_i, x_1) & \cdots & (x_n, x_1) \\
(x_1, x_2) & \cdots & (x_i, x_2) & \cdots & (x_n, x_2) \\
\vdots & & \vdots & & \vdots \\
(x_1, x_n) & \cdots & (x_i, x_n) & \cdots & (x_n, x_n)
\end{vmatrix}
}, \quad i = 1, 2, \cdots, n.
\end{align*}
\end{theorem}
\begin{remark}
\refexa{example:泛函分析--最小二乘法}中的三个问题本质上等价于这个定理.
\end{remark}
\begin{proof}
设 $x_0 = \sum\limits_{i=1}^{n} \lambda_i x_i$ 为所求元素，则由\refcor{corollary:泛函分析--推论289293}知其为 $x$ 在 $\mathscr{M}$ 上的正交投影当且仅当 $x - x_0 \perp \mathscr{M}$，即对任意 $j = 1, 2, \cdots, n$，有 $(x - x_0, x_j) = 0$.

将 $x_0 = \sum\limits_{i=1}^{n} \lambda_i x_i$ 代入内积条件，得
\begin{align*}
\left( x - \sum\limits_{i=1}^{n} \lambda_i x_i, x_j \right) = 0, \quad j = 1, 2, \cdots, n,
\end{align*}
展开即得
\begin{align}
\sum\limits_{i=1}^{n} \lambda_i (x_i, x_j) = (x, x_j), \quad j = 1, 2, \cdots, n. \label{eq:1_6_34}
\end{align}

由于 $x_1, x_2, \cdots, x_n$ 线性无关，Gram 矩阵 $((x_i, x_j))_{n \times n}$ 可逆，故方程组 \eqref{eq:1_6_34} 存在唯一解 $(\lambda_1, \lambda_2, \cdots, \lambda_n) \in \mathbb{R}^n$，即 $x_0$ 存在唯一.由Crammer法则可得
\begin{align*}
\lambda_i = \frac{
\begin{vmatrix}
(x_1, x_1) & \cdots & (x, x_1) & \cdots & (x_n, x_1) \\
(x_1, x_2) & \cdots & (x, x_2) & \cdots & (x_n, x_2) \\
\vdots & & \vdots & & \vdots \\
(x_1, x_n) & \cdots & (x, x_n) & \cdots & (x_n, x_n)
\end{vmatrix}
}{
\begin{vmatrix}
(x_1, x_1) & \cdots & (x_i, x_1) & \cdots & (x_n, x_1) \\
(x_1, x_2) & \cdots & (x_i, x_2) & \cdots & (x_n, x_2) \\
\vdots & & \vdots & & \vdots \\
(x_1, x_n) & \cdots & (x_i, x_n) & \cdots & (x_n, x_n)
\end{vmatrix}
}, \quad i = 1, 2, \cdots, n.
\end{align*}

\end{proof}


\begin{proposition}[极化恒等式]
设 $a$ 是复线性空间 $\mathscr{X}$ 上的共轭双线性函数, $q$ 是由 $a$ 诱导的二次型, 求证: 对 $\forall x,y\in\mathscr{X}$ 有
\begin{align*}
a(x,y)=\frac{1}{4}\left[q(x+y)-q(x-y)+\mathrm{i}q(x+\mathrm{i}y)-\mathrm{i}q(x-\mathrm{i}y)\right].
\end{align*}
\end{proposition}

\begin{example}
求证: 在 $C[a,b]$ 中不可能引进一种内积 $(\cdot,\cdot)$, 使其满足
\begin{align*}
(f,f)^\frac{1}{2}=\max_{a\leqslant x\leqslant b}|f(x)|\quad(\forall f\in C[a,b]).
\end{align*}
\end{example}

\begin{example}
在 $L^2[0,T]$ 中, 求证: 函数
\begin{align*}
x\mapsto\left|\int_0^T\mathrm{e}^{-(T-\tau)}x(\tau)\mathrm{d}\tau\right|\quad(\forall x\in L^2[0,T])
\end{align*}
在单位球面上达到最大值, 并求出此最大值和达到最大值的元素 $x$.
\end{example}
\begin{note}
提示 利用 Cauchy-Schwarz 不等式及其取等号的条件.
\end{note}

\begin{example}
设 $M,N$ 是内积空间中的两个子集, 求证:
\begin{align*}
M\subseteq N\implies N^\perp\subseteq M^\perp.
\end{align*}
\end{example}

\begin{example}
设 $M$ 是 Hilbert 空间 $\mathscr{X}$ 的子集, 求证:
\begin{align*}
(M^\perp)^\perp=\overline{\text{span}M}.
\end{align*}
\end{example}

\begin{example}
在 $L^2[-1,1]$ 中, 问偶函数集的正交补是什么? 证明你的结论.
\end{example}

\begin{example}
在 $L^2[a,b]$ 中, 考察函数集 $S=\{\mathrm{e}^{2\pi inx}\}$.
\begin{enumerate}[(1)]
\item 若 $|b-a|\leqslant 1$, 求证: $S^\perp=\{\theta\}$.
\item 若 $|b-a|>1$, 求证: $S^\perp\neq\{\theta\}$.
\end{enumerate}
\end{example}

\begin{example}
设 $\mathscr{X}$ 表示闭单位圆上的解析函数全体, 内积定义为
\begin{align*}
(f,g)=\frac{1}{i}\oint_{|z|=1}\frac{f(z)\overline{g(z)}}{z}\mathrm{d}z\quad(\forall f,g\in\mathscr{X}).
\end{align*}
求证: $\{z^n/\sqrt{2\pi}\}$ 是一组正交规范集.
\end{example}

\begin{example}
设 $\{e_n\}_{n=1}^\infty,\{f_n\}_{n=1}^\infty$ 是 Hilbert 空间 $\mathscr{X}$ 中的两个正交规范集, 满足条件
\begin{align*}
\sum\limits_{n=1}^\infty\|e_n-f_n\|^2<1.
\end{align*}
求证: $\{e_n\}$ 和 $\{f_n\}$ 两者中一个完备蕴含另一个完备.
\end{example}

\begin{example}
设 $\mathscr{X}$ 是 Hilbert 空间, $\mathscr{X}_0$ 是 $\mathscr{X}$ 的闭线性子空间, $\{e_n\},\{f_n\}$ 分别是 $\mathscr{X}_0$ 和 $\mathscr{X}_0^\perp$ 的正交规范基. 求证: $\{e_n\}\cup\{f_n\}$ 是 $\mathscr{X}$ 的正交规范基.
\end{example}

\begin{example}
设 $H^2(D)$ 是按例定义的内积空间.
\begin{enumerate}[(1)]
\item 如果 $u(z)$ 的 Taylor 展开式是 $u(z)=\sum\limits_{k=0}^\infty b_k z^k$, 求证:
\begin{align*}
\sum\limits_{k=0}^\infty\frac{|b_k|^2}{1+k}<\infty.
\end{align*}
\item 设 $u(z),v(z)\in H^2(D)$, 并且
\begin{align*}
u(z)=\sum\limits_{k=0}^\infty a_k z^k,\quad v(z)=\sum\limits_{k=0}^\infty b_k z^k,
\end{align*}
求证:
\begin{align*}
(u,v)=\pi\sum\limits_{k=0}^\infty\frac{a_k\overline{b_k}}{k+1}.
\end{align*}
\item 设 $u(z)\in H^2(D)$, 求证:
\begin{align*}
|u(z)|\leqslant\frac{\|u\|}{\sqrt{\pi}(1-|z|)}\quad(\forall|z|<1).
\end{align*}
\item 验证 $H^2(D)$ 是 Hilbert 空间.
\end{enumerate}
\end{example}

\begin{example}
设 $\mathscr{X}$ 是内积空间, $\{e_n\}$ 是 $\mathscr{X}$ 中的正交规范集, 求证:
\begin{align*}
\left|\sum\limits_{n=1}^\infty(x,e_n)\overline{(y,e_n)}\right|\leqslant\|x\|\cdot\|y\|\quad(\forall x,y\in\mathscr{X}).
\end{align*}
\end{example}

\begin{example}
设 $\mathscr{X}$ 是一个内积空间, $\forall x_0\in\mathscr{X},\forall r>0$, 令
\begin{align*}
C=\{x\in\mathscr{X}\mid\|x-x_0\|\leqslant r\}.
\end{align*}
\begin{enumerate}[(1)]
\item 求证: $C$ 是 $\mathscr{X}$ 中的闭凸集.
\item $\forall x\in\mathscr{X}$, 令
\begin{align*}
y=
\begin{cases}
x_0+r(x-x_0)/\|x-x_0\|, & \text{当}\ x\notin C, \\
x, & \text{当}\ x\in C.
\end{cases}
\end{align*}
求证: $y$ 是 $x$ 在 $C$ 中的最佳逼近元.
\end{enumerate}
\end{example}

\begin{example}
求 $(a_0,a_1,a_2)\in\mathbb{R}^3$, 使得
\begin{align*}
\int_0^1|\mathrm{e}^t-a_0-a_1t-a_2t^2|^2\mathrm{d}t
\end{align*}
取最小值.
\end{example}

\begin{example}
设 $f(x)\in C^2[a,b]$, 满足边界条件
\begin{align*}
f(a)=f(b)=0,\quad f'(a)=1,\quad f'(b)=0.
\end{align*}
求证:
\begin{align*}
\int_a^b|f''(x)|^2\mathrm{d}x\geqslant\frac{4}{b-a}.
\end{align*}
\end{example}

\begin{proposition}[变分不等式]
设 $\mathscr{X}$ 是一个 Hilbert 空间, $a(x,y)$ 是 $\mathscr{X}$ 上的共轭对称的双线性函数, $\exists M>0,\delta>0$, 使得
\begin{align*}
\delta\|x\|^2\leqslant a(x,x)\leqslant M\|x\|^2\quad(\forall x\in\mathscr{X}).
\end{align*}
又设 $u_0\in\mathscr{X}$, $C$ 是 $\mathscr{X}$ 上的一个闭凸子集. 求证: 函数
\begin{align*}
x\mapsto a(x,x)-\text{Re}(u_0,x)
\end{align*}
在 $C$ 上达到最小值, 并且达到最小值的点 $x_0$ 唯一, 且满足
\begin{align*}
\text{Re}\left[2a(x_0,x-x_0)-(u_0,x-x_0)\right]\geqslant 0\quad(\forall x\in C).
\end{align*}
\end{proposition}









\end{document}